\subsection{Systems - Steuerung des Verhaltens der beteiligten Komponenten}\label{sec:systems}

In einem Frame werden die über die \texttt{GameLoop} registrierten Systeme in einer bestimmten Reihenfolge aufgerufen.
Systeme iterieren hierbei nicht nur über einzelne Komponenten, sondern können - je nach Aufgabe - auch über eine Komponenten-Komposition iterieren, da beispielsweise eine Skalierungskomponente in unterschiedlichen Kontexten angewendet wird (UI, aktives NPC- oder Spielerobjekt).\footnote{
    Eine Komposition von Komponenten wird ausgedrückt durch das Vorhandensein einer Menge von verschiedenen Komponententypen.
    Unity bezieht sich darauf in ähnlicher Weise über die bereits erwähnten \textit{Archetypes}- / \textit{Chunks}~\cite{UnityChunks}.
}

\subsubsection{Tags statt Vererbung}

\textit{helios} drückt die Rolle eines Systems \textit{nicht} - entgegen der Implementierung in früheren Versionen - über Vererbung aus:

\begin{verbatim}
    // veraltet: die Rolle eines Systems wird über
    // Vererbung ausgedrückt
    class DamageOnCollisionSystem : public System {}
\end{verbatim}

\noindent
Stattdessen erfolgt die Kennzeichnung über \textit{Tags}, die zur Kompilierungszeit über \textit{Concepts} überprüft werden (siehe Listing~\ref{lst:system_inheritance}).
Damit geht die über Vererbung vermittelte Rollensemantik verloren, um zum Ausdruck zu bringen, dass die Rolle eines Systems nicht an einen bestimmten Basistypen gebunden ist - was das Paradigma der \textit{composition over inheritance} bekräftigen und Entwickler ermutigen soll, schlanke, leichtgewichtige \textit{Systeme} zu bevorzugen, worauf in den vorherigen Abschnitten mit dem Entwurfsziel von helios schon Stellung bezogen wurde.
Die hierbei genutzten \texttt{EngineRoleTags} werden auch in anderen Kernkonzepten von \textit{helios} verwendet, etwa bei \textit{Managern} und \textit{CommandHandlern} (siehe Abschnitte~\ref{sec:runtime_manager} und~\ref{sec:runtime_messaging}).

\begin{lstlisting}[style=c++style,
    caption={Auszug aus dem Quellcode der beteiligten Klassen und Concepts zur Erstellung von Systemen in \textit{helios}. Systeme müssen von keinem bestimmten Basistyp erben, dafür aber eine \texttt{update}-Methode besitzen, die einen \texttt{UpdateContext} als Argument erhält. Die Rolle eines Systems wird über Tags ausgedrückt, die über Concepts zur Kompilierungszeit überprüft werden.},
    label=lst:system_inheritance]

    // concept
    template<class T>
    concept IsSystemLike = requires(T& t, UpdateContext& updateContext) {
        {t.update(updateContext) } -> std::same_as<void>;
    } && HasTag<T, tags::SystemRole>;

    // umsetzende Klasse
    class DamageOnCollisionSystem {
        using EngineRoleTag = SystemRole;
        void update(UpdateContext& updateContext) {...}
    }

    // Compile-Time guard
    template<typename T, typename... Args>
    requires IsSystemLike<T>
    Pass& addSystem(Args&&... args) {
        systemRegistry_.template add<T>(
            std::forward<Args>(args)...
        );

        return *this;
    }

\end{lstlisting}

\subsubsection{Registrierung und Ausführung}

In dem Beispiel wurde die Implementierung eines \texttt{Pass} skizziert, der Teil der \texttt{GameLoop} ist, und das Concept \texttt{IsSystemLike} als Zulassungskriterium (\textit{Gatekeeping}) für die Registrierung von Systemen nutzt.
Das Concept stellt sicher, dass die übergebene Klasse die entsprechende \texttt{update}-Methode besitzt; der \texttt{Pass} ruft das System dann an der vorgesehenen Stelle innerhalb eines Frames auf.\\

Beim Aufruf erstellt die in Listing~\ref{lst:system_loop} gezeigte \texttt{update()}-Methode einen \textit{View} auf die erforderlichen Komponentenkonstellationen.
Die Logik ist ein Implementierungsdetail und kann über eine benötigte Aktualisierung der Komponenten hinausgehen und beispielsweise Nachrichten an nachgelagerte Systeme in demselben oder in nachfolgenden Frames senden.

\begin{lstlisting}[style=c++style,
    caption={Skizze einer \texttt{update()}-Implementierung, die einen View über aussschließlich aktive \textit{GameObjects} erstellt, bei denen die aufgeführten Komponenten auch \textit{enabled} sind. Auf die Rolle von \texttt{UpdateContext} wird in Abschnitt~\ref{sec:runtime_gameloop} näher eingegangen.},
    label=lst:system_loop]

    void update(UpdateContext& updateContext) {
        for (auto [go, cc, tsc, active] : updateContext.view<
                    ChaseComponent,
                    TranslationStateComponent,
                    Active
                >().whereEnabled()) {

            // ... Verhalten ...
        }
    }
\end{lstlisting}


